% ============================================================================
% Fichier  : partie5/P05_C16_csirt_cert.tex
% Titre    : Le CSIRT et le CERT
% Auteur   : MINKA MI NGUIDJOÏ Thierry Emmanuel
% ============================================================================
% Sources : ISO/IEC 27035 ; modèles de CSIRT/CERT ; ANTIC. Opération MAYO.
% ============================================================================

\chapter{Le CSIRT et le CERT}
\label{chap:csirt}
\chaptermark{Le CSIRT et le CERT}

\epigraph{%
    « En crise, ce qui n'est pas écrit n'a pas été décidé. »%
}{Principe du cours}

\niveauChapitre{Fondamental}{Chapitre~\ref{chap:cycle}.}

\begin{boxobjectifs}
\begin{itemize}
    \item Définir le CSIRT et le CERT, et leur différence.
    \item Décrire les rôles d'une cellule de réponse.
    \item Comprendre la main courante et la communication de crise.
    \item Situer la coopération avec l'ANTIC et les tiers.
    \item Dimensionner une cellule pour une PME.
\end{itemize}
\end{boxobjectifs}

\bigskip

% ============================================================================
\section*{Question directrice : qui décide, qui agit, qui parle pendant l'incident de LiZbiZ~?}
% ============================================================================

Sans organisation, un incident se gère dans la confusion~: tout le monde
agit, personne ne décide, et rien n'est tracé. Ce chapitre présente la
cellule de réponse, ses rôles, sa main courante et ses relations
extérieures.

\medskip
\textbf{Plan du chapitre.}
\begin{enumerate}
    \item CSIRT et CERT (section~\ref{sec:c16-definir}).
    \item Les rôles de la cellule (section~\ref{sec:c16-roles}).
    \item La main courante (section~\ref{sec:c16-maincourante}).
    \item Coopérer avec l'extérieur (section~\ref{sec:c16-cooperer}).
\end{enumerate}

\separateur

% ============================================================================
\section{CSIRT et CERT}
\label{sec:c16-definir}
% ============================================================================

\begin{boxdef}
\termecle{CSIRT} (\textit{Computer Security Incident Response Team})~: équipe
chargée de recevoir, analyser et traiter les incidents de sécurité d'un
périmètre donné. \termecle{CERT} est une appellation voisine, protégée~; tout
CERT est un CSIRT. L'\sigle{ANTIC} opère le CERT national camerounais.
\end{boxdef}

% ============================================================================
\section{Les Rôles de la Cellule}
\label{sec:c16-roles}
% ============================================================================

\begin{tableaucours}{L{3.6cm}L{8.6cm}}{%
    \tableauHeader{Rôle} &
    \tableauHeader{Responsabilité}
}
Responsable de la cellule & Décide, arbitre, rend compte à la direction \\
Analystes & Investiguent, contiennent, éradiquent \\
Référent technique métier & Connaît les systèmes atteints \\
Communication & Messages internes et externes, maîtrisés \\
Juridique & Obligations, signalements, ANTIC, assurance \\
Main courante & Consigne chaque décision et chaque action \\
\end{tableaucours}
\vspace{-0.8em}
\begin{center}{\small\itshape Tableau~16.1. Les rôles d'une cellule de réponse}\end{center}

% ============================================================================
\section{La Main Courante}
\label{sec:c16-maincourante}
% ============================================================================

\begin{boxprocedure}
\textbf{Tenir la main courante}
\begin{enumerate}
    \item Horodater chaque entrée.
    \item Noter la décision, qui l'a prise, sur quelle base.
    \item Noter l'action, qui l'a faite, son résultat.
    \item Ne rien reconstruire après coup~: on écrit sur le moment.
    \item Conserver la main courante au dossier (preuve et RETEX).
\end{enumerate}
\end{boxprocedure}

\begin{boxobjet}{Main courante de l'opération MAYO}
\textbf{Nature~:} journal horodaté des décisions et actions du 6~mars.
\textbf{Source~:} tenue en direct par le secrétaire de cellule.
\textbf{Pertinence~:} établit qui a décidé quoi, et quand.
\textbf{Fiabilité~:} entrées horodatées, non modifiées. \textbf{Opposabilité~:}
preuve de la diligence de LiZbiZ.
\end{boxobjet}

% ============================================================================
\section{Coopérer avec l'Extérieur}
\label{sec:c16-cooperer}
% ============================================================================

\begin{boxscenario}{Opération MAYO. La cellule de LiZbiZ}
\textbf{Constatation.} LiZbiZ n'avait pas de CSIRT. Le 6~mars, une cellule est
constituée en urgence~: M.~KOUNOU (technique), Mme~MOUAFO (responsable),
NYONG SÉCURITÉ (analystes), un juriste, un chargé de communication. L'ANTIC
est informée dans la journée.

\medskip
\textbf{Analyse} (confiance élevée). L'improvisation a coûté du temps.
Après l'incident, LiZbiZ formalise un CSIRT minimal et des playbooks~: la
prochaine cellule se réunira en minutes, pas en heures.
\end{boxscenario}

\begin{boxmetier}{Responsable de la réponse à incident}
\textbf{Désigner avant la crise.} Le responsable constitue la cellule et
nomme les rôles à froid, avec suppléants et coordonnées. En crise, on active
une liste, on ne cherche pas qui appeler.
\end{boxmetier}

\separateur

\begin{boxresume}
\textbf{Ce qu'il faut retenir de ce chapitre~:}
\begin{itemize}
    \item \textbf{CSIRT} (générique) et \textbf{CERT} (appellation
          protégée)~; l'ANTIC opère le CERT national.
    \item \textbf{Les rôles}~: responsable, analystes, métier, communication,
          juridique, main courante.
    \item \textbf{La main courante} trace décisions et actions, en direct.
    \item \textbf{On coopère} avec l'ANTIC et les tiers.
    \item \textbf{On désigne} la cellule avant la crise.
\end{itemize}
\end{boxresume}

\begin{boxinfo}
\textbf{Points de contrôle de l'élève.} Je sais~: distinguer CSIRT et CERT~;
citer les rôles~; tenir une main courante~; situer la coopération~;
dimensionner une cellule.
\end{boxinfo}

\bigskip

\begin{boxexo}[title={\ding{45}\quad Exercice~16.1. CSIRT ou CERT \niveaubase}]
Quelle est la différence entre CSIRT et CERT~?
\end{boxexo}

\begin{boxexo}[title={\ding{45}\quad Exercice~16.2. Rôles \niveaubase}]
Citez quatre rôles d'une cellule de réponse.
\end{boxexo}

\begin{boxexo}[title={\ding{45}\quad Exercice~16.3. Main courante \niveaubase}]
Que consigne-t-on dans une main courante~?
\end{boxexo}

\begin{boxexo}[title={\ding{45}\quad Exercice~16.4. Écrit \niveaumoyen}]
Pourquoi « ce qui n'est pas écrit n'a pas été décidé »~?
\end{boxexo}

\begin{boxexo}[title={\ding{45}\quad Exercice~16.5. ANTIC \niveaumoyen}]
Quand et pourquoi informer l'ANTIC~?
\end{boxexo}

\begin{boxexo}[title={\ding{45}\quad Exercice~16.6. MAYO \niveaumoyen}]
Qu'a coûté à LiZbiZ l'absence de CSIRT préétabli~?
\end{boxexo}

\begin{boxexo}[title={\ding{45}\quad Exercice~16.7. Communication \niveaumoyen}]
Pourquoi un rôle communication dédié dans la cellule~?
\end{boxexo}

\begin{boxexo}[title={\ding{45}\quad Exercice~16.8. Opposabilité \niveaumoyen}]
En quoi la main courante est-elle une preuve de diligence~?
\end{boxexo}

\begin{boxexo}[title={\ding{45}\quad Exercice~16.9. Charte de CSIRT \niveauexpert}]
Rédigez la charte du CSIRT minimal de LiZbiZ~: périmètre, rôles, suppléants,
coordonnées.
\end{boxexo}

\begin{boxexo}[title={\ding{45}\quad Exercice~16.10. Modèle de main courante \niveauexpert}]
Concevez le modèle de main courante de LiZbiZ.
\end{boxexo}

\bigskip

\begin{boxinfo}
\textbf{Pour aller plus loin.}
\begin{itemize}
    \item ISO/IEC 27035~\cite{iso27035}~; NIST SP~800-61r2~\cite{nist80061r2}.
\end{itemize}
\end{boxinfo}

\begin{boxinfo}
\textbf{Transition.} Pour anticiper, la cellule s'appuie sur le renseignement
sur la menace~: le Chapitre~\ref{chap:renseignement} le traite.
\end{boxinfo}
